\documentclass[letterpaper,12pt]{article}
\usepackage[T1]{fontenc}
\usepackage[utf8]{inputenc}
\usepackage[spanish]{babel}
\usepackage{framed}
\usepackage{amsfonts,setspace}
\usepackage{amsmath}
\usepackage{amssymb, amsmath, amsthm}
\usepackage{comment}
\usepackage[most]{tcolorbox}
\usepackage{amssymb}
\usepackage{dsfont}
\usepackage{anysize}
\usepackage{pdfpages}
\usepackage{multicol}
\usepackage{enumerate}
\usepackage{graphicx}
\usepackage{float}
\usepackage[dvipsnames]{xcolor}
\usepackage[left=1.5cm,top=0.9cm,right=1.5cm, bottom=1.4cm]{geometry}
\usepackage{fancyhdr}
\pagestyle{fancy}
\definecolor{blue(pigment)}{rgb}{0.2, 0.2, 0.6}
% Page
\oddsidemargin -0.4in
\textwidth 7in
\topmargin -0.6in
\textheight 9.1in
\parindent 0em
\parskip 1ex

%Teoremas, Lemas, etc.
\theoremstyle{plain}
\newtheorem{teo}{Teorema}
\newtheorem{lem}{Lemma}
\newtheorem{prop}{Proposici\'on}
\newtheorem{cor}{Corolario}
\theoremstyle{definition}
\newtheorem{defi}{Definici\'on}
\newtheorem{eje}{Ejemplo}
\newtheorem{obs}{Observaci\'on}
% fin macros

% Comandos más usados
% Conjuntos
\newcommand{\N}{\mathbb{N}}
\newcommand{\Z}{\mathbb{Z}}
\newcommand{\Q}{\mathbb{Q}}
\newcommand{\R}{\mathbb{R}}
\newcommand{\C}{\mathbb{C}}
\newcommand{\RR}{\mathbb{R} \times \mathbb{R}}
% Flechas
\newcommand{\ssi}{\Longleftrightarrow}
\newcommand{\imp}{\Longrightarrow}
\newcommand{\pmi}{\Longleftarrow}
% Trigonométricas
\newcommand{\sen}{\text{sen}}

\usepackage{versions}
\includeversion{solution}

\let\epsilon\varepsilon

\begin{document}
\def\Pauta{1}
%%%%%ESCRIBIR 1 si quieren ver la Prueba con Pauta%%%%
%%%%%ESCRIBIR 0 si quieren ver sólo la Ayudantía%%%%
	
	\fancyhead[L]{\itshape{Escuela de Ingeniería}}
	\fancyhead[R]{\itshape{Universidad de O'Higgins}}
	
	\begin{minipage}{11.5 cm}
         \vspace{-6mm}
		\begin{flushleft}
			\vspace{1mm}
			\textbf{ING1002-2 Cálculo Diferencial e Integral}\\
			\textbf{Profesor:}  Gonzalo Flores G.  \\
			\textbf{Ayudante:}  Juan I. Riquelme \& Cristian Acevedo R. \\
		\end{flushleft}
	\end{minipage}
	
	\begin{picture}(2,3)
		\put(430,9.5){\includegraphics[scale=0.22]{Imagenes/UOH.png}}
	\end{picture}
	\vspace{-6mm}
	\begin{center}
		\LARGE \bf{Ayudantía 10: Primitivas}
	\end{center}
    \vspace{-8.8mm}
    \begin{center}
		9 de octubre de 2025
	\end{center}

\begin{enumerate}[\textbf{P\arabic{enumi}.}]
\item Aplicar un cambio de variable para calcular las siguientes primitivas:
\begin{enumerate}
    \item $\displaystyle \int \frac{x}{\sqrt{x^2 - a^2}} \, dx$.

    \color{blue}
    \textbf{Sol:} Resolver con cambio de variable $u=x^2-a^2$.
    \color{black}

\item $\displaystyle
\int x e^{x^2-1}\,dx
$

\color{blue}    
\textbf{Sol:} Hacemos el cambio de variable \(u=x^2-1\), entonces
\[
du=2x\,dx.
\]
Obtenemos:
\[
\int x e^{x^2-1}\,dx
= \frac{1}{2}\int e^{u}\,du
= \frac{1}{2}e^{u}+C
= \frac{1}{2}e^{x^2-1}+C.
\]
    \color{black}
    
    
    \item $\displaystyle \int e^x \sqrt{1 + e^x} \, dx$.

\color{blue}
    \textbf{Sol:} \[
u=1+e^x \quad\Rightarrow\quad du=e^x\,dx.
\]
\[
\int e^x \sqrt{1+e^x}\,dx
= \int \sqrt{u}\,du
= \frac{2}{3}u^{3/2}+C.
\]
Volviendo a \(x\):
\[
\displaystyle \int e^x \sqrt{1+e^x}\,dx
= \frac{2}{3}\left(1+e^x\right)^{3/2}+C
\]
    \color{black}
    
    \item $\displaystyle \int x\sqrt{9-x^2}\, dx$.

    \color{blue}
\textbf{Sol:} 
Usamos el cambio de variable \(v=9-x^2\), con
\[
dv=-2x\,dx,
\]
con lo que
\[
\int x\sqrt{9-x^2}\,dx
= -\frac{1}{2}\int \sqrt{v}\,dv
= -\frac{1}{2}\cdot \frac{2}{3}v^{3/2}+C
= -\frac{1}{3}v^{3/2}+C.
\]
Volviendo a la variable \(x\):
\[
\int x\sqrt{9-x^2}\,dx
= -\frac{1}{3}(9-x^2)^{3/2}+C.
\]

    \color{black}
    
\end{enumerate}
\item Sean $f, g, h$ funciones tales que $f(x) = g'(x) + h'(x)g(x)$. Usando la definición de primitiva muestre que
\[
\int f(x)e^{h(x)}dx = e^{h(x)}g(x) + c.
\]
\color{blue}
    \textbf{Sol:} Dado que nos dan el valor de la primitiva
\[
F(x)=e^{h(x)}g(x)+c.
\]
\[
F'(x)=e^{h(x)}g'(x)+e^{h(x)}h'(x)g(x)
     =e^{h(x)}\bigl(g'(x)+h'(x)g(x)\bigr).
\]
dado que
\[
f(x)=g'(x)+h'(x)g(x),
\]
\[
F'(x)=f(x)e^{h(x)}.
\]
Por la definición de primitiva, se concluye que
\[
\int f(x)e^{h(x)}\,dx = e^{h(x)}g(x)+c.
\]
    \color{black}

\item Sean $f, g : [-1, 1] \to \mathbb{R}$ las funciones dadas por $f(x) = \sqrt{1 - x^2}$ y $g(x) = \sqrt{3} - \sqrt{1 - x^2}$.
\begin{enumerate}
    \item Calcular el área encerrada entre ambas curvas y las rectas $x = -1$ y $x = 1$.

\color{blue}
    \textbf{Sol:} propuesto
    \color{black}
    
    \item Determinar el volumen del sólido generado por la rotación de la región encerrada por el eje $OX$ y la curva $h(x) = \min\{f(x), g(x)\}$, en torno a $OX$.

\color{blue}
    \textbf{Sol:} propuesto 
    \color{black}
    
\end{enumerate}

\item Hallar el área de la región encerrada entre las parábolas $y = \frac{x^2}{3}$ e $y = 4 - \frac{2}{3}x^2$.

\color{blue}
    \textbf{Sol:} Hallamos primero los puntos de intersección para hallar en que intervalos una función está por arriba de la otra
\[
\frac{x^2}{3}=4-\frac{2}{3}x^2.
\]
\[
x^2=12-2x^2,
\qquad
3x^2=12,
\qquad
x^2=4,
\qquad
x=\pm2.
\]

En el intervalo $[-2,2]$ la parábola superior es
\(
y=4-\frac{2}{3}x^2
\)
y la inferior es
\(
y=\frac{x^2}{3}.
\), basta con sacar un valor del intervalo y reemplazar.
\[
A=\int_{-2}^{2}
\left[\left(4-\frac{2}{3}x^2\right)-\frac{x^2}{3}\right]dx
=\int_{-2}^{2}\left(4-x^2\right)dx.
\]

Calculando la integral,
\[
A=\left[4x-\frac{x^3}{3}\right]_{-2}^{2}
=\left(8-\frac{8}{3}\right)-\left(-8+\frac{8}{3}\right)
=\frac{32}{3}.
\]
    \color{black}


\item Sea $f : \mathbb{R} \to \mathbb{R}_{+}$ derivable y tal que $\displaystyle \int f(x)dx = f(x)$.
\begin{enumerate}
    \item Muestre que $\frac{f'(x)}{f(x)} = 1$ y deduzca que $\displaystyle \int \frac{f'(x)}{f(x)}dx = x + c$.

    \color{blue}
    \textbf{Sol:} Como
\[
\int f(x)\,dx=f(x),
\]
por definición de primitiva derivando a ambos lados, tenemos
\[
f(x)=f'(x).
\]
Como $f(x)>0$ para todo $x$, podemos dividir y obtenemos
\[
\frac{f'(x)}{f(x)}=1.
\]
\[
\int \frac{f'(x)}{f(x)}\,dx=\int 1\,dx=x+c.
\]

    \color{black}
    \item Concluya que $f(x) = e^{x+c}$.

    \color{blue}
    \textbf{Sol:} De lo anterior tenemos
\[
\int \frac{f'(x)}{f(x)}\,dx=x+c.
\]
Haciendo un cambio de variable $f(x)=u \rightarrow f'(x)dx=du$, entonces
\[
\int \frac{f'(x)}{f(x)}\,dx=\int\frac{du}{u}=\ln(u)=\ln f(x)
\]
por lo tanto
\[
\ln f(x)=x+c.
\]
Exponenciando ambos lados
\[
f(x)=e^{x+c}
\]


    \color{black}
\end{enumerate}

\item Considere las funciones $F, G : [0, 1] \to \mathbb{R}$ dadas por
    \[
    F(x) = \int_0^x t e^{-t} dt, \quad G(x) = \int_0^x \cos(\pi x) e^{-t^2} dt.
    \]
    Sea además $H : [0, 1] \to \mathbb{R}$ la función dada por
    \[
    H(x) = F(x) - 2G'(x).
    \]
    \begin{enumerate}
        \item Justifique que $H$ es Riemann integrable en $[0, 1]$.

        \color{blue}
    \textbf{Sol:} Por TFC
\[
F'(x)=x e^{-x}.
\]
\[
G(x)=\cos(\pi x)\int_0^x e^{-t^2}\,dt.
\]
\[
G'(x)
=-\pi\sen(\pi x)\int_0^x e^{-t^2}\,dt
+\cos(\pi x)e^{-x^2}.
\]
\[
H(x)=F(x)-2G'(x)
=\int_0^x t e^{-t}\,dt
+2\pi\sen(\pi x)\int_0^x e^{-t^2}\,dt
-2\cos(\pi x)e^{-x^2}.
\]
Cada término de esta expresión es continuo en $[0,1]$, luego $H$ es continua en $[0,1]$. Como toda función continua en un intervalo cerrado y acotado es Riemann integrable, se concluye que
\[
H \text{ es Riemann integrable en } [0,1].
\]

    \color{black}
        
        \item Demuestre que
        \[
        H(1) = 1.
        \]
        \textbf{Indicación:} puede usar que $\displaystyle \int_0^1 t e^{-t} dt = -2e^{-1} + 1$.

        \color{blue}
    \textbf{Sol:} Evaluamos en $x=1$:
\[
H(1) = (-2 e^{-1} + 1) + 0 + 2 e^{-1} = 1.
\]

    \color{black}
    \end{enumerate}


\color{blue}
    \textbf{Propuesto}
\color{black}

\item \textbf{[Propuesto].-} Calcule las siguientes primitivas usando cambio de variable

    
\begin{enumerate}
    \item $\displaystyle \int \frac{1}{x} \ln(x) dx$

    \color{blue}
    \textbf{Sol:} Sea \( u = \ln(x) \), entonces derivando, \( du = \frac{1}{x} dx \), reemplazando
\[
\int u \, du = \frac{u^2}{2} + C = \frac{(\ln(x))^2}{2} + C
\]
Por lo tanto:
\[
\int \frac{\ln(x)}{x} \, dx = \frac{(\ln(x))^2}{2} + C
\]
    \color{black}
\item $\displaystyle \int \frac{\cos(x) e^{\frac{3}{\sen(x)}}}{\sen^2(x)} dx$

 \color{blue}
    \textbf{Sol:} Realizamos el siguiente cambio de variable:
\[
u = \frac{3}{\sen(x)} \ \Rightarrow \
du = \frac{d}{dx} \left( \frac{3}{\sen(x)} \right) 
= -\frac{3 \cos(x)}{\sen^2(x)} \, dx
\ \Rightarrow  - \frac{\sen^2(x)}{3 \cos(x)} \, du = dx
\]
Sustituimos
\[
\int \frac{\cos(x) \cdot e^u}{\sen^2(x)} \cdot dx 
= \int \frac{\cos(x) \cdot e^u}{\sen^2(x)} \cdot 
\left( - \frac{\sen^2(x)}{3 \cos(x)} \, du \right)
\]
\[
= - \frac{1}{3} \int e^u \, du
= - \frac{1}{3} e^u + C
\]
Desaciendo el cambio de variable
\[
= - \frac{1}{3} e^{\frac{3}{\sen(x)}} + C
\]
\color{black}
\item $\displaystyle \int \frac{e^{\ln(x)}}{x} dx$

 \color{blue}
    \textbf{Sol:} Queda propuesto el desarrollo con cambio de variable, pero decir que se puede resolver facilmente de esta manera\[
\int \frac{e^{\ln(x)}}{x} \, dx = \int \frac{x}{x} \, dx = \int 1 \, dx = x + C \quad\quad \textcolor{blue}{\text{(Recordando que}\  e^{\ln(x)}=x)}
\]
\color{black}
\item $\displaystyle\int \frac{x^2 + 1}{x - 1} dx$

 \color{blue}
    \textbf{Sol:} Dividimos el polinomio:
\[
\frac{x^2 + 1}{x - 1} = x + 1 + \frac{2}{x - 1}
\]
Por lo tanto,
\[
\int \frac{x^2 + 1}{x - 1} \, dx = \int (x + 1) \, dx + \int \frac{2}{x - 1} \, dx = \frac{x^2}{2} + x + 2 \ln|x - 1| + C
\]
    \color{black}
\end{enumerate}


\end{enumerate}


\end{document}